\section{The continuation problem and the main result}
\label{ag:sec:problem}

The manuscript's Gamma hierarchy has two distinct kinds of questions.
The first is algebraic: which differential monomials survive a specified
integration-by-parts reduction, and with which rational coefficients?
The second is analytic: how do their integrals, spectral derivatives,
and normalized primitives reduce to special functions? Neither question
should be replaced by an assertion of independence of numerical constants.

The immediate source for this continuation is \emph{Exact Harmonic-Tail
Reconstruction and Variational Gamma Identities} \cite{HTG}. Its
Section~4 determines the complete normal projection onto $u_0,u_1,u_2$.
Its final research question~5 asks for the next projection, including
$u_3^2$ with arbitrary powers of $u_1,u_2$. We solve the subfamily with
no additional $u_2$ factor, and simultaneously solve its analogue for
every derivative order $r\ge2$. In particular, the result is not only
a table at the next weight.

Let $\D u_j=u_{j+1}$ in $\Q[u_0,u_1,\ldots]$. Write $\NF$ for the
canonical projection onto monomials which are either powers of $u_0$
or have their largest positive index repeated at least twice. Define
\begin{equation}\label{ag:eq:hierarchy-intro}
 A_0=1,\qquad A_{m+1}=u_0A_m-\D A_m,\qquad N_m=\NF(A_m).
\end{equation}
The decomposition and all conventions are proved or specified in the
next section. Under $u_j=\psi^{(j)}(x)$,
\[
 A_m(x)=(-1)^m\Gamma(x)\frac{\mathrm d^m}{\mathrm dx^m}\frac1{\Gamma(x)}.
\]

\begin{theorem}[All derivative-square coefficients]\label{ag:thm:main}
Let $r\ge2$ and $a,k\ge0$, and put $m=a+2k+2r+2$. Then
\begin{equation}\label{ag:eq:coefficient}
 \boxed{
 [u_0^a u_1^k u_r^2]N_m=
 \begin{cases}
 \displaystyle
 \frac{(-1)^r m!}
 {a!\,2^{k+r+1}k!(r+1)!\prod_{j=0}^{r}(k+2j+1)},&k\text{ even},\\[7pt]
 0,&k\text{ odd}.
 \end{cases}}
\end{equation}
Equivalently, with $y=u_1t^2/2$,
\begin{equation}\label{ag:eq:normal-sector}
 \sum_{m\ge0}N_m\frac{t^m}{m!}
 =e^{tu_0}\left\{f_0(y)+\sum_{r\ge2}u_r^2t^{2r+2}g_r(y)
                          +R(t;u_1,u_2,\ldots)\right\},
\end{equation}
where every monomial of $R$ contains at least three factors from
$\{u_2,u_3,\ldots\}$, and
\begin{align}
 f_0(y)&=\frac12\left(\cosh y+\frac{\sinh y}{y}-y\Shi(y)\right),
       \label{ag:eq:f0-intro}\\
 g_r(y)&=\frac{(-1)^r}{(r+1)!}\piy^{\,r+1}\cosh y
 =\frac{(-1)^r}{(2r+2)!}\,{}_0F_1\left(;r+\frac32;\frac{y^2}{4}\right).
       \label{ag:eq:gr-intro}
\end{align}
Here $\Shi(y)=\int_0^y\sinh v\dd v/v$, and
$\piy=(2y)^{-1}\partial_y$ acts on even power series. All apparent
singularities at $y=0$ are removable.
\end{theorem}

The phrase ``quadratic sector'' refers to the number of factors having
index at least two. It does not mean the total factor degree, and it does
not discard arbitrary powers of $u_0$ or $u_1$. Formula
\eqref{ag:eq:normal-sector} is an identity of formal series, not an
approximation to a Gamma ratio on an interval. Its proof is in
Section~\ref{ag:sec:polarization}.

\subsection{The overlap and the advance}

At $k=0$, \eqref{ag:eq:coefficient} is
$(-1)^r\binom{m}{2r+2}$; the inspected predecessors already contain
these pure-square coefficients. At $r=2$ it agrees with the complete
$u_2^2$ family in \cite[Section~4]{HTG}. Those cases are checks and
inherited specializations, not additional discoveries claimed here.
The advance is the single all-$r$, all-$k$ formula and its affine-symbol
proof, together with the symmetric-profile and analytic consequences.

For example, the coefficients of $u_1^{2\ell}u_r^2$ at their corresponding
weights $4\ell+2r+2$ begin as follows. These are not coefficients at
one fixed weight.
\begin{center}
\begin{tabular}{@{}rrrrr@{}}
\toprule
$r$ & $\ell=0$ & $\ell=1$ & $\ell=2$ & $\ell=3$\\
\midrule
2&1&90&15015&4176900\\
3&$-1$&$-165$&$-40950$&$-15261750$\\
4&1&273&96390&47006190\\
5&$-1$&$-420$&$-203490$&$-127193220$\\
6&1&612&395010&310810500\\
\bottomrule
\end{tabular}
\end{center}
Multiplication by $u_0^a$ multiplies the exponential-series coefficient
by $1/a!$, as displayed in the theorem.

The analytic part is connected by the beta representation of $g_r$,
not by an unproved interchange of the entire Gamma normal-form series
with a singular integral. It yields finite Hurwitz and Lerch endpoint
formulas, every generalized Stieltjes order, and ordinary integral
identities such as
\begin{equation}\label{ag:eq:trigamma-preview}
 \int_1^3(x-1)^2(3-x)^2\psi'(x)\dd x=\frac{6\zeta(3)}{\pi^2}.
\end{equation}
Individual polynomially weighted special-function integrals can have
classical antecedents; our claim is the proved structural derivation
and its explicitly scoped extension, not universal historical priority.
